\chapter*{Cathetan Sumber lan Notasi}
\addcontentsline{toc}{chapter}{Cathetan Sumber lan Notasi}

Ing tuladha relasi tatanan, sumber nganggo $I$ tanpa
ngenali jenenge luwih dhisik. Ing kono, $I$ kudu diwaca
minangka relasi identitas ing wilangan natural,
yaiku pasangan-pasangan ing diagonal kang wis diterangake.

Ing definisi pang wit, sumber nulis $z \in X \setminus B$,
nanging himpunan simpul wit wis dijenengi $A$ lan $X$
durung ditemtokake. Ing definisi iku, wacanen $X$ minangka $A$.
Maksimalitas banjur ateges saben simpul ing njaba $B$
ora bisa dibandhingake karo paling ora siji simpul ing $B$.

Notasi $R^+$ nduweni teges lokal kang beda:
ing perangan tatanan, tegese tutupan refleksif;
ing perangan operasi relasi, tegese tutupan transitif.
Gunakna definisi ing perangan kang lagi diwaca.

Andharan ekstensionalitas ing wiwitan kudu diwaca
bebarengan karo perangan Paradoks Russell:
ekstensionalitas njamin mung ana siji himpunan
kang cocog yen himpunan kasebut pancen ana.

Ing tuladha fungsi oyod kuadrat, sumber nyebut oyod kang dipilih
minangka positif, nanging domain wilangan natural kalebu nol.
Ing nol, wacanen iki minangka oyod kuadrat utama kang ora negatif.
Iki temuan audit sumber OLFUN-002 lan wis dibenerake ing teks terjemahan.
Ing tuladha himpunan pandhisik lan panerus, masukan dijenengi $n$,
nanging banjur ditulis $x$. Ing kono, $x$ nuduhake masukan $n$ kang padha.
Iki temuan OLFUN-003; teks terjemahan saiki nggunakake $x$ kanthi ajeg.

Proposisi yen saben injeksi duwe invers kiwa mbutuhake katrangan
tumrap domain kosong. Bukti ing sumber milih $a \in A$, mula
lumaku yen $A$ ora kosong. Yen $A$ lan $B$ loro-lorone kosong,
fungsi kosong dadi inverse dhewe. Nanging yen $A$ kosong lan $B$
ora kosong, injeksi kosong ora duwe invers kiwa saka $B$ menyang $A$.
Mula, wacanen proposisi iku kanthi sarat $A$ ora kosong utawa $B$ kosong.
Iki temuan OLFUN-001; awak proposisi terjemahan nganggo sarat prasaja
$A$ ora kosong lan cathetan jejer menehi sarat umum kang pas.
Panganggone Aksioma Pilihan ing bukti invers tengen wis diterangake
ing cathetan sikil sumber.

Ing andharan graf fungsi, sumber nyebut $R_f\subseteq A\times B$
minangka relasi "ing $A\times B$". Miturut definisi buku iki,
relasi ing himpunan kasebut malah dadi subhimpunan
$(A\times B)^2$. Teks terjemahan wis mbenerake iki dadi relasi
antarane $A$ lan $B$ kang kaemot ing $A\times B$ (OLFUN-004).

Watesan fungsi ing $C$ mung nyaring koordinat masukan, mula grafe
$R_f\cap(C\times B)$. Watesan relasi ing bagean sadurunge yaiku
$R\cap C^2$ lan nyaring loro koordinat. Rumus fungsi ing sumber bener;
pratelan yen loro gagasan persis padha wis diwatesi dadi analogi
kanthi bedane dijlentrehake (OLFUN-005).

Cathetan iki tambahan panyunting ing terjemahan mesin.
Rumus ing sumber selaras tetep dijaga supaya
bisa dibandhingake langsung karo sumber Inggris.
